\section{Benoit Replication: Reproduction Details}\label{app:benoit-replication}

Section~\ref{sec:benoit-headline} summarizes the Benoit replication on
C-TreePO. This appendix specifies the corpus, scoring protocol, empirical
comparisons, and full per-cell grid behind the main-text summary.

\subsection{Corpus, Targets, and Expert Benchmarks}

The main comparison uses the same Manifesto Project platform benchmark
as \citet{BenoitEtAl2025}: $235$ party election platforms from
parliamentary democracies, filtered to election years for which the
Benoit--Laver and CHES expert surveys provide party-position estimates.
Benoit et al.'s broader paper also reports $23$ Kl\"uver coalition
agreements where applicable, but the headline C-TreePO comparison is
organized around the Manifesto Project platform rows. Raw text is taken
from the Manifesto Corpus 2025a release \citep{ManifestoProject2025a};
where the release text contained promotional footer truncation, we used
the full official Manifesto Project text.

The prediction target is an expert-survey mean. Benoit et al.'s
replication archive provides one harmonized mean party-position estimate
for each (manifesto, dimension) pair. Using that archive rather than
recomputing from raw CHES plus Benoit--Laver avoids rescaling,
party-year-join, and missing-value discrepancies. As a check, we
recompute Benoit's published Figure~1 correlations from their released
LLM-score and expert-mean tables and reproduce their numbers to within
$0.005$ on every dimension.

Benoit et al.'s Table~3 reference is a split-expert agreement
benchmark, not a sample-specific limit for our held-out split. In the
Benoit--Laver replication data, experts within each country are randomly
divided into two groups. Each group produces a mean party-position
estimate on the same issue-position scale, and the two vectors of mean
estimates are correlated. The economic value used below is
$r = 0.880$; the macro value across the six policy dimensions is
$0.873$.

\subsection{Pipeline Protocol}

Every row in Table~\ref{tab:benoit-headline} corresponds to a single
pipeline configuration. Each configuration has three components:

\begin{itemize}[leftmargin=1.5em]
    \item \textbf{Leaf summarizer.} Each text span is summarized under
          a dimension-specific preservation rubric. The prompt targets
          Benoit's $300$--$400$ word summary length but allows longer
          outputs when aggressive compression would drop policy-relevant
          evidence.
    \item \textbf{Merge summarizer.} Adjacent child summaries are merged
          under the same rubric. The merge prompt has the same
          information-preservation target as the leaf prompt and is
          allowed to carry both child summaries forward when that is the
          least lossy state.
    \item \textbf{Scorer.} The root summary is scored with the Benoit
          rubric for dimension $d$ and mapped to the seven-point policy
          scale, with a missing response recorded only when the model
          cannot produce a valid score.
\end{itemize}

The six scoring contexts are the exact prompts released in Benoit's
replication materials, so the scorer sees the same rubric text as the
GPT-4o scorer in their pipeline. This matters: an earlier paraphrased
rubric shifted macro correlation by $+0.008$ on held-out summaries when
we swapped back to the exact released text.

\paragraph{Chunking.} Chunks are produced by
paragraph-boundary splitting inside a character budget. Typical MP
manifestos are $20$--$60$K characters; at $c = 8\mathrm{K}$ this yields
$3$--$8$ leaves, and at $c = 24\mathrm{K}$ it yields $1$--$3$ leaves.

\paragraph{Model.} All runs use
\textsc{Gemma-4-31B-IT-NVFP4} served from a local GPU node with a
$12\mathrm{K}$-token context cap. The context cap is more than
sufficient for the summaries used here ($\sim\!3\mathrm{K}$ tokens)
and keeps the inference budget fixed across settings.

\paragraph{Scoring schedule.} The joint variant of
Section~\ref{sec:benoit-headline} makes $235$ summaries $\times$ $6$
dimensions $= 1{,}410$ scoring calls. Each call's SYSTEM prompt is
one of six Benoit rubrics ($\sim\!1.1$K tokens each). Iterating the
pipeline in manifesto-then-dimension order
interleaves all six rubrics across consecutive calls. The joint variant
therefore uses a two-pass evaluation schedule: first generate all root
summaries, then score all manifestos for one dimension before moving to
the next. The $\sim\!235$ scoring calls for each dimension then share
the same rubric prefix. Saved prompt-prefix work is
$\sim\!1.1\text{K tokens} \times (N-1) \times 6
\approx 1.5$M tokens per full-corpus run.

\subsection{Full Per-Cell Results}

Table~\ref{tab:benoit-headline} in the main text reports the macro
summary. The complete per-dimension $\times$ chunk-size $\times$
configuration grid includes:

\begin{itemize}[leftmargin=1.5em]
    \item Per-dimension C-Tree at $c\in\{2,4,8,16,24,32,64\}\text{K}$.
    \item Joint C-Tree at $c\in\{8,16,24,32,64\}\text{K}$.
    \item \textbf{Flat baseline} (no chunking, no summarization, raw
          text truncated to $c$ characters and scored directly): macro
          $0.786$--$0.827$ across truncation sizes. This is the direct
          long-context measurement control: it tells us how much of the
          expert-survey signal is visible before the text is distributed
          over leaves.
    \item \textbf{Concat baseline} (chunking and leaf summaries but no
          merges --- leaf summaries are concatenated and scored):
          distributes the measurement across leaves and then recombines
          them by simple concatenation. Macro $0.827$--$0.837$ at
          $c\in\{8\mathrm{K},16\mathrm{K}\}$ shows that the expert-survey
          signal survives the first local compression step.
    \item \textbf{Decoding sensitivity.} Additional low-temperature and
          medium-temperature runs, with and without three-sample
          aggregation, leave the ranking of the main configurations
          essentially unchanged. At $c = 8\mathrm{K}$ the per-dimension
          C-Tree macro lies in $[0.832,0.841]$ across these variants,
          while the joint C-Tree lies in $[0.800,0.808]$.
    \item \textbf{Gemma-3 exact-model replication}. Running the
          scorer on Benoit's own anonymized summaries with the
          Gemma-3-27B BF16 model used in their Table~6 produces macro
          $0.723$--$0.731$ across rubric and prompt-format
          variations. The gap to their reported $0.793$ is not
          explained by rubric wording ($+0.008$ delta) or output-format
          wrapping ($-0.003$ delta); the residual is attributed
          elsewhere to model-weight revision drift between the
          Benoit vintage and our local checkpoint.
    \item \textbf{Scorer-only baselines}. Running only the scorer
          stage of our pipeline on Benoit's anonymized summaries
          (i.e.\ no tree) gives a macro $r$ in the $0.73$--$0.75$
          band. The jump to $0.82$--$0.84$ in
          Table~\ref{tab:benoit-headline} is attributable to the
          summarization step, not to the scorer alone.
\end{itemize}

\subsection{Ablations on the Mergeability Story}

Two ablations load directly onto Proposition~\ref{prop:mergeable-reduction}'s
empirical prediction.

\paragraph{Chunk-invariance across a $16\times$ chunk sweep.} We ran
the per-dimension C-Tree at $c\in\{4,8,16,24,32,64\}\mathrm{K}$ on the
same manifestos; the macro spread across chunk sizes is $0.027$, with
the five non-decentralization dimensions staying in the same
high-correlation regime while decentralization is the visible exception
discussed next. This is the empirical
analog of \textbf{Prop.~\ref{prop:mergeable-reduction}}: under local
validity the root-level score should be chunk-schedule-invariant, and
we observe the predicted behavior to within sampling noise on five
dimensions.

\paragraph{The dimension where chunk-invariance breaks is the
dimension the published literature flags.} Decentralization's
Pearson $r$ monotonically \emph{improves} as leaves shrink: $0.489$
at $c = 64\mathrm{K}$, $0.543$ at $c = 8\mathrm{K}$, $0.584$ at
$c = 2\mathrm{K}$. \citet{BenoitEtAl2025} flag decentralization as a
boundary case in their Figures~2--3, where manifesto quasi-sentence
coders and expert-survey respondents disagree about what
          resolves a party's position; their own Table~3 split-expert
          reference on decentralization is $0.78$, well below the
          $\ge 0.88$ references for the other five dimensions. The smaller-leaf regime in
C-TreePO forces the tree to carry the disambiguating evidence across
merges rather than summarizing it away at the leaf, and the
correlation improves as a result. This is the Markov counter of
Section~\ref{sec:markov-walkthrough} manifesting in real data: when
the oracle cares about between-leaf interactions, finer leaves help.

\paragraph{Flat and concat controls.} The flat baseline ($c$-character
truncation, no tree) reports macro $0.786$--$0.827$ depending on
truncation, with the strongest row at a $48\mathrm{K}$-character
truncation. The concat baseline (leaf summaries, no recursive merges)
reports $0.827$--$0.837$, including $0.837$ at both $32\mathrm{K}$ and
$8\mathrm{K}$ leaves. These rows are measurement-distribution controls.
The flat row measures the signal available to a direct read; the concat
row measures the signal after the text has already been spread across
leaf summaries. Their strength is useful because it shows that the
Benoit benchmark has a strong recoverable signal. The C-Tree rows then
show that the same expert-survey signal survives the fully local architecture: macro
$0.829$ at $8\mathrm{K}$ and $0.842$ at $4\mathrm{K}$, with stable
correlations across the chunk sweep. The target fact is invariance as
the measurement grid moves from direct reads to many local leaves. Once
that invariance holds, the researcher can sample, supervise, and audit a
larger set of smaller measurement points while preserving the
root-level target.

\subsection{Additional Robustness Diagnostics}

\paragraph{Manifesto length.} Binning platforms by raw manifesto length
shows that the headline correlations are not driven only by short
documents. Among manifestos longer than $100\mathrm{K}$ characters, the
per-dimension C-Tree reports Pearson $r=0.89$ on economic policy,
$0.85$ on social policy, $0.87$ on immigration, $0.88$ on European
integration, and $0.83$ on environment. Decentralization remains the
weak dimension at $0.47$, matching the broader pattern in
Table~\ref{tab:benoit-headline} rather than creating a new long-document
failure mode. The joint C-Tree shows the same qualitative pattern in the
longest bucket: the five non-decentralization dimensions remain in the
$0.85$--$0.91$ range.

\paragraph{Election era.} Splitting the full per-dimension pipeline by
election year gives the same picture. At $c = 24\mathrm{K}$, economic
correlations are $0.899$ for 1989--1999, $0.859$ for 2000--2009, and
$0.905$ for 2010--2019; social correlations are $0.847$, $0.791$, and
$0.865$; environment correlations are $0.802$, $0.759$, and $0.830$.
Immigration and European integration have no 1989--1999 cells in this
benchmark split, but remain high in the 2000s and 2010s
($0.801$--$0.889$ for immigration and $0.862$--$0.912$ for European
integration). Decentralization again stays lower, from $0.393$ to
$0.497$. Thus the main result is not concentrated in one recent era of
manifesto rhetoric.

\paragraph{Secondary Manifesto Project coded-score target.} Benoit et
al. also discuss Manifesto Project quasi-sentence counts transformed
into coded logit scores. We treat those as a construct-validity check,
not as the target benchmark for this section, because they measure a
different object from expert-survey party-position means. Against the
coded-score target, the tree, flat, concat, and joint variants produce
macro correlations in the $0.337$--$0.397$ range, and the European
integration correlations are strongly negative. This confirms that the
expert-survey target and the coded-count target should not be treated as
interchangeable validation labels.

\subsection{Compute Cost Accounting}

The compute cost decomposes cleanly by pipeline variant and LLM call
type.

\begin{itemize}[leftmargin=1.5em]
    \item \textbf{Benoit proprietary ensemble.} Per manifesto--dimension:
          $3 \text{ LLMs} \times 2 \text{ shot-settings} \times 3 \text{ summaries} = 18$
          calls (nine summarization calls, nine scoring calls). Corpus
          total: $235\times 6\times 18 = 25{,}380$ calls, routed
          through GPT-4o, Claude 3.5 Sonnet, and Gemini 1.5 Pro
          frontier APIs.
    \item \textbf{Joint C-Tree} (our default for the full sweep).
          Per manifesto: $2\lceil n_{\text{chars}}/c\rceil - 1$
          summarize/merge calls plus $6$ scoring calls. At
          $c = 8\mathrm{K}$ and a typical $40$K--$48$K-character
          manifesto, that is $15$--$17$ calls. Corpus total is on the
          order of $2{,}500$ LLM calls on one local model.
    \item \textbf{Per-dimension C-Tree.} Six trees per manifesto;
          $\approx 60$--$70$ LLM calls per manifesto at
          $c = 8\mathrm{K}$; corpus total $\approx 15{,}000$ calls on
          one local model.
\end{itemize}

At frontier API list prices ($\sim\!\$0.01$--$\$0.05$ per $1$K
tokens) the Benoit ensemble is on the order of a few hundred dollars
per corpus run; the same run on a local GPU node for $\sim\!8$ hours
of wall-clock is a rental cost roughly two orders of magnitude smaller
on a dollar-to-dollar basis, and rate-limited vs.\ bursty for
iteration. No claim is made that this is a head-to-head cost
comparison --- it depends on hardware ownership, amortization,
and the cost of calling frontier APIs at scale --- but the
qualitative story (one local open-weight model vs.\ three frontier
APIs $\times$ six passes each) is clean.

\subsection{Reproduction}

The full sweep regenerates the tables in Section~\ref{sec:benoit-headline}
from per-manifesto predictions in four protocol steps: evaluate the
open-weight model; run the per-dimension pipeline over the six Benoit
dimensions and chunk sizes
$c \in \{4, 8, 16, 24, 32, 64\}\mathrm{K}$; run the joint two-pass
pipeline over $c \in \{8, 16, 24, 32, 64\}\mathrm{K}$; and aggregate the
resulting manifesto-level scores against Benoit's expert means. The
aggregation unit is always the manifesto--dimension prediction, so the
reported correlations do not depend on the order in which predictions
are produced.

\subsection{Oracle-Grounded Training Path}

The headline Benoit replication uses a prompt-only summarizer and
scorer with no training against internal-node oracles. Expert
quasi-sentence codings are available for $2{,}157$ platforms
($\sim\!2.27$M coded spans). These annotations define a span oracle:
for any leaf or internal-node span, aggregate the quasi-sentence codes
that fall inside that span and return the corresponding rubric-indexed
score. This gives ground-truth oracle values at every leaf (whatever
span the leaf covers) and at every internal node (the union-span of its
two children), so each of the three local laws becomes an explicit
training signal:

\begin{itemize}[leftmargin=1.5em]
    \item \textbf{C1 (leaf sufficiency):} the head applied to each
          leaf state must predict the leaf-span score to within the
          coder's resolution.
    \item \textbf{C3 (per-merge consistency):} the head applied to each
          internal-node state must predict the ordered union-span score.
          An optional order-swap probe on the root merge can be reported
          when the scoring rubric treats quasi-sentence evidence as
          permutation-invariant, but this is a stress test rather than
          the definition of C3.
    \item \textbf{C2 (idempotence):} re-summarizing an already-produced
          summary must not shift the predicted score.
\end{itemize}

This supervision path mirrors the operator-learning setup used in the
Markov and HLL experiments in
Sections~\ref{sec:markov-walkthrough}--\ref{sec:hll-parity}; it uses
an embedding-space FNO operator as $g$ and a shared neural head as
$f$, with strict C1/C3 supervision at every node and a C2 probe on
the root merge. Section~\ref{sec:manifesto-qsentence} reports the
realized experiments on this path, with full grids in
Appendices~\ref{app:qsentence-leaf-ladder}
and~\ref{app:gold-seeding}.

\subsection[Alternating f/g Optimization on the Economic Policy Dimension]{Alternating $f/g$ Optimization on the Economic Policy
Dimension}\label{app:rile-fg-ladder}

This subsection provides the full details behind
Section~\ref{sec:rile-fg-ladder}'s main-body summary of the alternating
$f/g$ ladder.

\paragraph{Methodology.} Each ladder step runs medium-budget
prompt-program search on one side of $(f, g)$ with the other side's
prompt held fixed at the previous iteration's output. Weights of the
underlying Gemma-4-31B-NVFP4 checkpoint are never updated; the optimized
object is the prompt program. Stage names use the compact label
$f^a g^b$: $a$ indexes the number of scorer-side prompt optimizations,
and $b$ indexes the number of summarizer-side prompt optimizations. The
notation is a stage label, not multiplication. Leaf size is reported in
\emph{tokens}; the character-based leaf column in
Table~\ref{tab:benoit-headline} is a distinct axis and should not be
compared at face value.

\paragraph{Two initialization lanes.} The pilot uses two starting
lanes:
\begin{itemize}[leftmargin=1.5em]
    \item \textbf{Benoit-init lane:} $g^0 = g^{\mathrm{Benoit}}$, meaning the
          archived GPT-4o summary from Benoit's replication materials.
          The plotted anchor is $f^1 g^0$ because the scorer has already
          gone through the scorer-only optimization reported in the
          previous subsections. The lane then asks what happens when the
          summarizer is updated against that scorer.
    \item \textbf{Raw-init lane:} $g^0$ is the
          own-Gemma-4 baseline summary generated under the scoring
          rubric; $f^0$ is the untuned scorer. Both sides of the
          pair come from the same model, so the first ladder
          step $(f^0 g^0 \to f^1 g^0)$ is the cleanest possible
          co-adaptation vignette.
\end{itemize}
When a later ladder stage reused the same evaluation as the preceding
stage, Table~\ref{tab:manifesto-fg-ladder} reports only the first stage
at which the distinct number appeared.

\paragraph{Heatmap figures.} Figure~\ref{fig:manifesto-fg-ladder-heatmaps}
shows external-Pearson heatmaps for both lanes, with rows indexing the
ladder stage and columns indexing leaf size in tokens.

\begin{figure}[t]
    \centering
    \begin{subfigure}[b]{0.48\textwidth}
        \centering
        \includegraphics[width=\linewidth]{manifesto_fg_ladder_benoit_init.png}
        \caption{Benoit-init lane ($g^0 = g^{\mathrm{Benoit}}$). Populated
        across leaves $256$--$8192$ tok.}
        \label{fig:manifesto-fg-ladder-benoit}
    \end{subfigure}
    \hfill
    \begin{subfigure}[b]{0.48\textwidth}
        \centering
        \includegraphics[width=\linewidth]{manifesto_fg_ladder_raw_init.png}
        \caption{Raw-init lane ($g^0$ = own-Gemma baseline). Leaves
        $256$--$8096$ tok populated at the $f^1 g^0$ and $f^1 g^1$
        stages; deeper stages fill when a summarizer update changed
        the evaluation, and are shown with a $=$ glyph where an
        evaluation was inherited from the preceding stage.}
        \label{fig:manifesto-fg-ladder-raw}
    \end{subfigure}
    \caption{Paired heatmaps for the two initialization lanes of the
    alternating $f/g$ ladder on the economic policy dimension
    (seven-point \citet{BenoitEtAl2025} scale, $n = 48$ held-out
    manifestos per cell). Both lanes use \textsc{Gemma-4-31B-IT-NVFP4}
    with medium-budget prompt-program search applied to one side of
    $(f, g)$ per step while the other is held fixed. Rows are ladder
    stages. In the label $f^a g^b$, $a$ counts scorer-side prompt
    optimizations and $b$ counts summarizer-side prompt optimizations;
    this is a stage label, not multiplication. Bottom to top, the full
    raw-init ladder is
    $f^0 g^0 \to f^1 g^0 \to f^1 g^1 \to f^2 g^1 \to f^2 g^2 \to
    f^3 g^2 \to f^3 g^3$; the Benoit-init lane omits $f^0 g^0$ because
    it starts from the scorer-optimized anchor $f^1 g^0$. Columns: leaf
    size in tokens.
    \textbf{Left panel of each lane:} external Pearson $r$ against the
    Benoit expert-survey mean. The bolded tick at $r = 0.880$ on the
    colorbar marks the economic split-expert reference from
    \citet{BenoitEtAl2025}'s Table~3: the correlation between mean
    party-position estimates from two random expert subsamples. Cells
    above the tick exceed that empirical reference on this held-out split;
    cells below the tick sit beneath it. The tick is not a theoretical
    limit. \textbf{Right panel of each lane:} internal-external
    Pearson gap (lower = better), defined as Pearson agreement with the
    teacher trace minus Pearson agreement with Benoit's expert means.
    The dotted line at $0$ marks parity between the internal and
    external correlations. Repeated values across stages in the
    raw-init panel ($f^1 g^1$ through $f^3 g^3$ at leaf $256$) show
    evaluations inherited from the most recent post-update
    stage, not separate re-runs.}\label{fig:manifesto-fg-ladder-heatmaps}
\end{figure}

\paragraph{Per-cell grid.} Table~\ref{tab:manifesto-fg-ladder} reports
the full ladder grid for both lanes. Columns cover external Pearson
against the expert mean, external MAE on the $1$--$7$ scale, the
internal $f$ Pearson (scorer applied to the realized summary vs. the
teacher trace for that same summary), and the internal-external Pearson gap.

\begin{table}[htbp]
    \centering
    \resizebox*{!}{0.88\textheight}{\resizebox{\linewidth}{!}{% Auto-generated by paper/ctreepo/scripts/render_ladder_table.py.
% Collapses inherited (unchanged) ladder stages into the first
% stage at which each distinct evaluation appeared.
\begin{tabular}{lrrrrrr}
\toprule
Lane & Leaf (tok) & Stage & Ext.\ $r$ & Ext.\ MAE & Int.\ $r$ & Int.-Ext.\ gap \\
\midrule
\multicolumn{7}{l}{\textit{Benoit-init lane: $g^0 = g^{\mathrm{Benoit}}$, $f^1$ = scorer-only optimized prompt}} \\
Benoit & 256 & $f^{1} g^{0}$ & $0.750$ & $1.657$ & --- & --- \\
Benoit & 256 & $f^{1} g^{1}$ & $0.824$ & $1.558$ & $0.985$ & $0.161$ \\
Benoit & 256 & $f^{2} g^{2}$ & $0.830$ & $1.555$ & $0.990$ & $0.160$ \\
Benoit & 512 & $f^{1} g^{1}$ & $0.861$ & $1.427$ & $0.982$ & $0.121$ \\
Benoit & 512 & $f^{2} g^{1}$ & $0.831$ & $1.513$ & $0.966$ & $0.134$ \\
Benoit & 512 & $f^{2} g^{2}$ & $0.834$ & $1.541$ & $0.953$ & $0.120$ \\
Benoit & 1024 & $f^{1} g^{1}$ & $0.885$ & $1.501$ & $0.974$ & $0.089$ \\
Benoit & 1024 & $f^{2} g^{2}$ & $0.874$ & $1.487$ & $0.969$ & $0.095$ \\
Benoit & 2048 & $f^{1} g^{1}$ & $0.874$ & $1.445$ & $0.968$ & $0.093$ \\
Benoit & 2048 & $f^{2} g^{2}$ & $0.886$ & $1.564$ & $0.977$ & $0.092$ \\
Benoit & 4096 & $f^{1} g^{1}$ & $0.865$ & $1.605$ & $0.989$ & $0.124$ \\
Benoit & 4096 & $f^{2} g^{1}$ & $\mathbf{0.886}$ & $1.543$ & $0.983$ & $0.097$ \\
Benoit & 4096 & $f^{2} g^{2}$ & $0.863$ & $1.626$ & $0.975$ & $0.111$ \\
Benoit & 8192 & $f^{1} g^{1}$ & $0.866$ & $1.543$ & $0.973$ & $0.107$ \\
Benoit & 8192 & $f^{2} g^{2}$ & $0.879$ & $1.501$ & $0.965$ & $\mathbf{0.086}$ \\
\midrule
\multicolumn{7}{l}{\textit{Raw-init lane: $g^0$ = own-Gemma baseline summary, $f^0$ = untuned scorer}} \\
Raw & 256 & $f^{0} g^{0}$ (anchor) & $0.826$ & $1.766$ & --- & --- \\
Raw & 256 & $f^{1} g^{0}$ & $0.873$ & $1.564$ & $0.976$ & $0.103$ \\
Raw & 256 & $f^{1} g^{1}$ & $\mathbf{0.909}$ & $1.543$ & $0.965$ & $\mathbf{0.056}$ \\
Raw & 512 & $f^{1} g^{0}$ & $0.861$ & $1.526$ & $0.970$ & $0.109$ \\
Raw & 512 & $f^{1} g^{1}$ & $0.871$ & $1.522$ & $0.970$ & $0.099$ \\
Raw & 512 & $f^{2} g^{2}$ & $0.868$ & $1.706$ & $0.967$ & $0.099$ \\
Raw & 512 & $f^{3} g^{3}$ & $0.862$ & $1.647$ & $0.970$ & $0.108$ \\
Raw & 1024 & $f^{1} g^{0}$ & $0.872$ & $1.626$ & $0.980$ & $0.108$ \\
Raw & 1024 & $f^{1} g^{1}$ & $0.853$ & $1.657$ & $0.969$ & $0.116$ \\
Raw & 1024 & $f^{2} g^{1}$ & $0.847$ & $1.654$ & $0.969$ & $0.122$ \\
Raw & 2048 & $f^{1} g^{0}$ & $0.888$ & $1.519$ & $0.976$ & $0.088$ \\
Raw & 2048 & $f^{1} g^{1}$ & $0.860$ & $1.519$ & $0.962$ & $0.102$ \\
Raw & 4096 & $f^{1} g^{0}$ & $0.885$ & $1.501$ & $0.973$ & $0.087$ \\
Raw & 4096 & $f^{1} g^{1}$ & $0.858$ & $1.505$ & $0.972$ & $0.114$ \\
Raw & 8096 & $f^{1} g^{0}$ & $0.868$ & $1.543$ & $0.975$ & $0.107$ \\
Raw & 8096 & $f^{1} g^{1}$ & $0.855$ & $1.730$ & $0.980$ & $0.125$ \\
Raw & 8096 & $f^{2} g^{1}$ & $0.853$ & $1.584$ & $0.971$ & $0.117$ \\
Raw & 8096 & $f^{2} g^{2}$ & $0.870$ & $1.577$ & $0.966$ & $0.096$ \\
Raw & 8096 & $f^{3} g^{3}$ & $0.867$ & $1.547$ & $0.971$ & $0.104$ \\
\bottomrule
\end{tabular}
}}
    \caption{Per-cell ladder grid on the economic policy dimension
    (seven-point scale).
    Stage labels use the notation defined in
    Figure~\ref{fig:manifesto-fg-ladder-heatmaps}.
    Both lanes collapse inherited stages (where a later stage reused
    the preceding evaluation) into the first stage at which each
    distinct evaluation appeared. Bold external-$r$ values are the
    best cell within each lane; bold internal-external gap values are
    the tightest (closest to parity) within each
    lane.}\label{tab:manifesto-fg-ladder}
\end{table}

\paragraph{Leaf-size robustness.} Collapsing
Table~\ref{tab:manifesto-fg-ladder} to the per-leaf best cell in the
Benoit-init lane gives $r = 0.830$ (leaf $256$), $0.861$ ($512$),
$0.885$ ($1024$), $0.886$ ($2048$), $0.886$ ($4096$), $0.879$ ($8192$).
From leaf $1024$ onward the realized correlation sits in a $0.007$-wide
band ($[0.879, 0.886]$) straddling the \citet{BenoitEtAl2025} Table~3
split-expert reference of $r = 0.880$ across an $8\times$ leaf-size
range; below $1024$ tok the fall-off is sharp ($-0.025$ at leaf $512$
and $-0.056$ at leaf $256$, measured relative to the plateau peak of
$0.886$). This is the same chunk-invariance
pattern Table~\ref{tab:benoit-headline} reports for the fixed-prompt
pipeline ($\pm 0.027$ macro across $16\times$ leaf variation), now
replicated under iterative training rather than fixed prompts, which
is the non-obvious result: medium-budget alternating prompt updates do not
introduce chunk-size sensitivity that the fixed-prompt version did
not already have, and they leave the realized operator near the
economic split-expert reference rather than clearly below it.

\paragraph{Internal-consistency panel.} The internal-external Pearson
gap reference is $0$: a gap of zero means the scorer's agreement with
the teacher trace is no larger than its agreement with Benoit's expert
means. Across every populated Benoit-init cell, the internal-$f$ Pearson
sits in $[0.953, 0.990]$ and the internal-external Pearson gap sits in
$[0.086, 0.161]$, with the minimum gap at leaf $8192$, $f^2 g^2$
($0.086$). Across the leaves where the Benoit-init Pearson plateau holds
($\geq 1024$ tok) the gap is $\leq 0.124$. The internal agreement
probe therefore stays tight across the ladder: alternating prompt
updates preserve high agreement with the teacher trace while the
external correlation tracks the expert benchmark. The raw-init lane
reaches the grid's tightest gap ($0.056$) at leaf $256$, $f^1 g^1$,
where the same cell also posts the grid's best external Pearson
($r = 0.909$); across the six leaves, the raw-init gap stays in
$[0.056, 0.125]$, matching the Benoit-init range on the leaves where
both lanes are populated.

\paragraph{Apples-to-apples comparison.} The relevant baseline for the
Benoit-init ladder is the scorer-only ablations block of
Table~\ref{tab:benoit-headline} (Benoit summaries held fixed, various
prompt optimizers on $f$ alone, economic $r$ in $[0.817, 0.837]$). The
best Benoit-init ladder cell reaches $r = 0.886$, a $+0.049$
improvement over the strongest scorer-only ablation ($0.837$) and
$+0.064$ over the few-shot optimized scorer ($0.822$). The
    ladder's gain is therefore localized to the summarizer-update half of
    the loop: re-optimizing $g$ against the current $f$ is what moves the
    system above the fixed-$g$ ablation range. We do not compare the ladder's numbers
to the per-dimension headline row of Table~\ref{tab:benoit-headline}
($0.939$ on economic at $c = 8\mathrm{K}$ chars) because the leaf axis
units differ and because the matched-leaf fixed-prompt baseline at the
same prompt-search budget has not been run; see the main-body caveats in
Section~\ref{sec:rile-fg-ladder}.

\paragraph{What remains.} Two pieces are flagged for follow-up: the
extension of the ladder to all six Benoit policy dimensions is scoped
but not yet launched; and a fixed-prompt baseline at matched token-leaf
sizes and the same prompt-search budget would let the ladder numbers
be compared directly against the chunk-invariance figures of
Table~\ref{tab:benoit-headline} rather than against the scorer-only
ablations block. Both are tracked as pending updates rather than as
claims in the present submission.
